\subsection{Related Work}\label{sec:fundamentals-relatedwork}

Three prior systems attempt in-context medical image segmentation, each
making a different choice on the three axes this thesis organizes around.

\paragraph{UniverSeg~\cite{butoiUniverSegUniversalMedical2023}} popularized
the approach with a CrossBlock architecture that densely cross-attends
target and context features at full resolution -- no coarse-to-fine
structure, and no synthetic-task training.

\paragraph{Iris~\cite{gaoShowSegmentUniversal2025}} encodes each context
pair into a compact per-class task embedding and fuses it into the target
through a pixel shuffle-unshuffle operation -- a single fusion step,
applied once, with no iterative refinement across scales.

\paragraph{Medverse~\cite{huMedverseUniversalModel2025a}} encodes target
and context with two separate U-Nets, fused via blockwise cross-attention
between the two streams at every stage, pooling queries and keys to block
resolution for efficiency (\autoref{sec:methodology-fusion}).

\paragraph{Hierarchical, coarse-to-fine approaches.} Deep Neural
Patchworks~\cite{reisertDeepNeuralPatchworks2022a} stacks levels of
increasing resolution for large-volume segmentation, but targets
single-model supervised segmentation rather than in-context, few-shot
prediction -- an architectural precedent for coarse-to-fine cascades
outside the ICL setting. Medverse applies a similar idea within the ICL
setting itself: it performs next-scale autoregressive segmentation,
inspired by next-scale image generation~\cite{tian2024visual}, where at
each step the previous step's (perturbed) prediction re-enters as an
additional context pair, processed by the same context U-Net -- reused
unchanged across scales. Every step still re-scans the \emph{entire}
volume with a sliding window -- only the input changes with scale, not
the spatial extent processed.

\begin{table}[t]
  \centering
  \small
  \caption{Positioning against prior in-context segmentation models on
  the three axes this thesis organizes around.}\label{tab:relatedwork}
  \begin{tabular}{lccc}
    \toprule
     & Fusion & Coarse-to-fine & Synthetic tasks \\
    \midrule
    UniverSeg & dense cross-attn. & no & no \\
    Iris & one-shot pixel shuffle & no & no \\
    Medverse & dual U-Net cross-attn. & full-volume re-scan & no \\
    \textbf{Ours} & dual attention (repeated) & region-restricted, register-carried & GMM repaint + texture \\
    \bottomrule
  \end{tabular}
\end{table}
